\documentclass[../../main.tex]{subfiles}
\graphicspath{{../../image/}{../../figure/}}
\begin{document}

\section{一般系统}\label{section:ch13-一般系统}

本节扼要讨论，对特征族仅为分段真非线性的系统，如何得到近似解 $U_h$ 沿网格曲线的总变差界。这些界将用\cref{section:ch13-glimm泛函}中对真非线性系统所用的步骤导出，只是度量未来波相互作用势的泛函须修改，因为此处波相互作用受\cref{chapter:ch09-容许波扇与riemann问题}（而非\cref{chapter:ch09-容许波扇与riemann问题}）支配。
于是，考虑菱形 $\Delta_s^r$，其入射波扇 $\alpha$ 与 $\beta$ 分别通过``西南''边与``东南''边进入，出射波扇为 $\varepsilon$，则\cref{chapter:ch09-容许波扇与riemann问题}给出
\begin{align}
|\varepsilon-(\alpha+\beta)|\leqslant c_{11}\mathscr{D}(\Delta_s^r),
\label{align:ch13-gen-int}
\end{align}
其中
\begin{align}
\mathscr{D}(\Delta_s^r)=\sum\theta|\gamma||\delta|.
\label{align:ch13-gen-D}
\end{align}
回忆上述求和遍历所有初等 $i$ 波（振幅 $\gamma$）与 $j$ 波（振幅 $\delta$）对，它们分别通过 $\Delta_s^r$ 的``西南''边与``东南''边进入。权重因子 $\theta$ 确定如下：若 $i<j$，$\theta=0$；若 $i>j$ 或 $i=j$ 且 $\gamma\delta<0$，$\theta=1$；最后，若 $i=j$ 且 $\gamma\delta>0$，$\theta$ 由\cref{chapter:ch09-容许波扇与riemann问题}给出。

如\cref{section:ch13-glimm泛函}中那样，对任意网格曲线 $I$ 相联系由\eqref{align:ch13-L-def}定义的泛函 $\mathscr{L}(I)$。设 $J$ 是如\cref{figure:ch13-网格曲线}所示的 $I$ 的直接后继，\eqref{align:ch13-gen-int}给出
\begin{align}
\mathscr{L}(J)\leqslant\mathscr{L}(I)+c_{11}\mathscr{D}(\Delta_s^r).
\label{align:ch13-gen-L-change}
\end{align}

\eqref{align:ch13-gen-L-change}所允许的 $\mathscr{L}$ 的增大，将由泛函 $\mathscr{Q}$ 的减小所抵消，$\mathscr{Q}$ 监测未来波相互作用势，此处定义为
\begin{align}
\mathscr{Q}(I)=\sum\theta|\zeta||\xi|.
\label{align:ch13-gen-Q}
\end{align}
上述求和遍历穿越网格曲线 $I$ 的所有初等 $i$ 波与 $j$ 波对，其振幅分别为 $\zeta$ 与 $\xi$。当 $i$ 波在 $j$ 波之左穿越 $I$ 时，若 $i<j$ 则 $\theta=0$，若 $i>j$ 则 $\theta=1$。当 $i=j$ 且 $\zeta\xi<0$ 时，$\theta=1$。最后，若 $i=j$ 且 $\zeta\xi>0$，$\theta$ 由\cref{chapter:ch09-容许波扇与riemann问题}确定；特别地，当左方波为速度 $\sigma_L$ 的 $i$ 激波、右方波为速度 $\sigma_R$ 的 $i$ 激波时由\cref{chapter:ch09-容许波扇与riemann问题}$_1$确定，
或当左方波为速度 $\sigma_L$ 的 $i$ 激波、右方波为连接 $U_R$ 与 $V_i(\tau_R;U_R)$ 的 $i$ 稀疏波时由\cref{chapter:ch09-容许波扇与riemann问题}$_2$确定，
或当左方波为连接 $U_L$ 与 $V_i(\tau_L;U_L)$ 的 $i$ 稀疏波、右方波为速度 $\sigma_R$ 的 $i$ 激波时由\cref{chapter:ch09-容许波扇与riemann问题}$_3$确定，
或当左方波为连接 $U_L$ 与 $V_i(\tau_L;U_L)$ 的稀疏波、右方波亦为连接 $U_R$ 与 $V_i(\tau_R;U_R)$ 的稀疏波时由\cref{chapter:ch09-容许波扇与riemann问题}$_4$确定。

目标是证明\eqref{align:ch13-Q-change}的类比，即若 $J$ 是如\cref{figure:ch13-网格曲线}所示的 $I$ 的直接后继，则
\begin{align}
\mathscr{Q}(J)-\mathscr{Q}(I)\leqslant[c_{12}\mathscr{L}(I)-1]\mathscr{D}(\Delta_s^r).
\label{align:ch13-gen-Q-change}
\end{align}

一旦建立\eqref{align:ch13-gen-Q-change}，便如\cref{section:ch13-glimm泛函}中那样，考虑由\eqref{align:ch13-G-def}定义的 Glimm 泛函 $\mathscr{G}$，并证明当 $\kappa$ 取得充分大、$\mathscr{L}(I)$ 很小时 $\mathscr{G}(J)\leqslant\mathscr{G}(I)$。这进而由\cref{section:ch13-总变差界}中所用论证给出所需估计\eqref{align:ch13-h-tv}与\eqref{align:ch13-h-lip}。

为验证\eqref{align:ch13-gen-Q-change}，我们重走\eqref{align:ch13-Q-change}证明中的步骤，并作必要调整。仍用``外围波''与``主体波''这两个术语，以区分同时穿越 $I$ 与 $J$ 的初等波与进入或离开 $\Delta_s^r$、从而仅穿越 $I$ 或仅穿越 $J$ 的初等波。

首先，两个入射波扇 $\alpha$ 与 $\beta$ 中主体波的相互作用向 $\mathscr{Q}(I)$ 贡献量 $\mathscr{D}(\Delta_s^r)$。相反，出射波扇 $\varepsilon$ 中主体波对不对 $\mathscr{Q}(J)$ 作贡献。

下一观察是，外围波对向 $\mathscr{Q}(I)$ 与 $\mathscr{Q}(J)$ 的贡献相等；因而它们对 $\mathscr{Q}(J)-\mathscr{Q}(I)$ 的净贡献为零。

剩下须考察外围波与主体波的配对。估计某个固定外围 $i$ 波（振幅 $\zeta$）与 $\alpha,\beta,\varepsilon$ 中初等 $j$ 波的配对对 $\mathscr{Q}(I)$ 与 $\mathscr{Q}(J)$ 的贡献。如真非线性情形，须区分如下情形：(i) $j>i$ 且外围 $i$ 波在 $(y_s^{r-1},t_s)$ 之左穿越 $I$；(ii) $j<i$ 且外围 $i$ 波在 $(y_s^{r+1},t_s)$ 之右穿越 $I$；(iii) $j>i$ 且外围 $i$ 波在 $(y_s^{r+1},t_s)$ 之右穿越 $I$；(iv) $j<i$ 且外围 $i$ 波在 $(y_s^{r-1},t_s)$ 之左穿越 $I$；(v) $j=i$，$\alpha_i\beta_i>0$ 且 $\zeta(\alpha_i+\beta_i)<0$；(vi) $j=i$，$\alpha_i\beta_i<0$ 且 $\zeta(\alpha_i+\beta_i)<0$；(vii) $j=i$，$\alpha_i\beta_i>0$ 且 $\zeta(\alpha_i+\beta_i)>0$；(viii) $j=i$，$\alpha_i\beta_i<0$ 且 $\zeta(\alpha_i+\beta_i)>0$。

在情形 (i)、(ii) 与 (iii) 中，外围 $i$ 波不与 $\alpha,\beta,\varepsilon$ 中任何 $j$ 波趋近；因而对 $\mathscr{Q}(I)$ 与 $\mathscr{Q}(J)$ 的贡献均为零。相反，在情形 (iv)、(v) 与 (vi) 中，对 $\mathscr{Q}(I)$ 与 $\mathscr{Q}(J)$ 的贡献分别为 $|\zeta|(|\alpha_j|+|\beta_j|)$ 与 $|\zeta||\varepsilon_j|$。

在情形 (v) 中，对 $\mathscr{Q}(I)$ 的贡献为 $|\zeta|(|\alpha_i|+|\beta_i|)$。对 $\mathscr{Q}(J)$ 的贡献依赖于 $\zeta\varepsilon_i$ 的符号，但在任何情况下不超过量 $|\zeta||\varepsilon_i|$。类似地，在情形 (vi) 中，对 $\mathscr{Q}(I)$ 的贡献至少为 $|\zeta|\max\{|\alpha_i|,|\beta_i|\}$，而对 $\mathscr{Q}(J)$ 的贡献至多为 $|\zeta||\varepsilon_i|$。

由上述及\eqref{align:ch13-gen-int}，属于情形 (i) 至 (vi) 之一的外围 $i$ 波与所有主体波配对，对 $\mathscr{Q}(J)-\mathscr{Q}(I)$ 的总贡献不超过量 $c_{11}|\zeta|\mathscr{D}(\Delta_s^r)$。

其余情形 (vii) 与 (viii) 需要更精细的处理。事实上，正是在此处真非线性系统与一般系统的区别显现出来。作为引导，考察\cref{chapter:ch09-容许波扇与riemann问题}证明中所考虑的特殊但有代表性的情形：
入射波扇 $\alpha$ 与 $\beta$ 各由单个 $i$ 激波组成，振幅分别为 $\gamma$ 与 $\delta$，速度分别为 $\sigma_L$ 与 $\sigma_R$，$\sigma_R\leqslant\sigma_L$。出射波扇 $\varepsilon$ 的第 $i$ 个波亦由单个 $i$ 激波组成，振幅 $\varepsilon_i$，速度 $\sigma$。为确定起见，进一步假设外围 $i$ 波同样是 $i$ 激波，振幅 $\zeta$，速度 $\sigma_0<\sigma_R$，在 $(y_s^{r+1},t_s)$ 之右穿越 $I$。按上述情形 (vii)，令 $\gamma,\delta,\varepsilon_i$ 与 $\zeta$ 均为正。则对 $\mathscr{Q}(J)-\mathscr{Q}(I)$ 的贡献为
\begin{align}
\zeta\Big\{(\sigma-\sigma_0)^+\varepsilon_i-(\sigma_L-\sigma_0)^+\gamma-(\sigma_R-\sigma_0)^+\delta\Big\},
\label{align:ch13-gen-edge}
\end{align}
由\cref{chapter:ch09-容许波扇与riemann问题}与\cref{chapter:ch09-容许波扇与riemann问题}，它为 $O(1)\zeta\theta\gamma\delta$。
一般情形中 $\alpha,\beta$ 为任意入射波扇的证明需要冗长的技术分析，但遵循相同模式，其中\cref{chapter:ch09-容许波扇与riemann问题}与\cref{chapter:ch09-容许波扇与riemann问题}扮演\cref{chapter:ch09-容许波扇与riemann问题}与\cref{chapter:ch09-容许波扇与riemann问题}的角色；
见\cref{section:ch13-注记}引用的文献。最终结论是：振幅为 $\zeta$ 的任意外围波与所有主体波配对，对 $\mathscr{Q}(J)-\mathscr{Q}(I)$ 的总贡献不超过量 $c_{12}|\zeta|\mathscr{D}(\Delta_s^r)$。因此，此类相互作用对 $\mathscr{Q}(J)-\mathscr{Q}(I)$ 的总贡献被 $c_{12}\mathscr{L}(I)\mathscr{D}(\Delta_s^r)$ 控制。这就建立了\eqref{align:ch13-gen-Q-change}，从而给出 $U_h$ 的总变差界。

在\cref{section:ch13-注记}引用的文献中表明，上述估计甚至可推广到更一般的严格双曲守恒律系统类，即可由特征族分段真非线性的系统``一致地''逼近的系统。这一更广的类例如涵盖\cref{chapter:ch07-一维双曲平衡律组}的等熵弹性动力学，
其中应力--应变曲线任意光滑、严格递增。

\end{document}